\section{Word Break}
\label{sec:rec-word-break}

\problemstatement{Given a string $s$ and a dictionary of words
$\textit{wordDict}$, determine whether $s$ can be segmented into a
sequence of one or more dictionary words placed consecutively (words may
be reused any number of times).}

\begin{patternbox}
Structurally, this is Palindrome Partitioning's
(Section~\ref{sec:rec-palindrome-partitioning}) ``loop over the next
piece's extent'' choice shape, combined with Sudoku's and Word Search's
boolean short-circuit-on-success control flow -- but with one crucial
difference that makes this section the deliberate bridge into the
dynamic programming chapters that follow: the sub-problem
``can $s[\textit{start}..]$ be segmented'' depends only on
$\textit{start}$, and the \textbf{same} value of $\textit{start}$ can be
reached through many different earlier choices of piece boundaries --
these are \textbf{overlapping subproblems}, the defining signature of a
problem where memoization will pay off enormously.
\end{patternbox}

\subsection*{Understanding the problem carefully}

From position $\textit{start}$, try every possible end position for the
next word: if $s[\textit{start}..\textit{end}]$ is in the dictionary,
recursively check whether the remainder, $s[\textit{end}+1..]$, can also
be fully segmented. If reaching the very end of the string is possible
through any sequence of such choices, $s$ is breakable. As in Sudoku
(Section~\ref{sec:rec-sudoku-solver}), the moment any choice leads to a
fully successful segmentation of the remainder, we can stop trying
further alternatives and propagate \texttt{true} immediately.

\begin{ideabox}[Loop Over Dictionary-Matching Prefixes, Short-Circuit on Success]
$\textsc{Solve}(\textit{start})$: if $\textit{start} = |s|$: return
\texttt{true} (the whole string has been consumed). Otherwise, for
$\textit{end}$ from $\textit{start}$ to $|s|-1$: if
$s[\textit{start}..\textit{end}]$ is in $\textit{wordDict}$ and
$\textsc{Solve}(\textit{end}+1)$ returns \texttt{true}: return
\texttt{true} immediately. If no choice of $\textit{end}$ leads to
success: return \texttt{false}.
\end{ideabox}

\subsection*{Why this recursion, run naively, repeats the same work exponentially often}

Consider a string with many ways to reach the same intermediate position
$\textit{start}$ -- for instance, $s=\texttt{"aaaaaaaaaaaaaaaaaaaaaaaab"}$
with dictionary $\{\texttt{"a"},\texttt{"aa"},\texttt{"aaa"}\}$: the
position right after consuming, say, the first $10$ characters can be
reached by choosing word lengths $1+1+\cdots$, or $2+2+\cdots$, or
$1+2+1+\cdots$, and many other combinations -- yet
$\textsc{Solve}(10)$ always asks exactly the same question (``can
$s[10..]$ be segmented'') regardless of \emph{which} combination of
earlier word choices got us to position $10$. A naive recursive
implementation recomputes the answer to $\textsc{Solve}(10)$ separately,
from scratch, every single time any path happens to reach it -- and the
number of distinct paths reaching a given position can grow
exponentially with the string's length, even though there are only
$|s|+1$ \emph{distinct} values $\textit{start}$ could ever take.

\subsection*{The bridge to dynamic programming}

This is precisely the situation dynamic programming exists to fix:
\textbf{memoization} caches the result of $\textsc{Solve}(\textit{start})$
the first time it is computed, in an array (or hash map) indexed by
$\textit{start}$, and every subsequent call with the same
$\textit{start}$ simply looks up the cached answer in $O(1)$ instead of
re-deriving it -- collapsing the exponential blowup down to $O(|s|)$
distinct sub-problems, each solved once. The recursion's
\emph{structure} (the recurrence relating $\textsc{Solve}(\textit{start})$
to $\textsc{Solve}(\textit{end}+1)$ for various $\textit{end}$) does not
change at all; only the bookkeeping around it does. This exact
transformation -- recognize overlapping subproblems in an already-correct
recursion, then cache results by the parameters that define each distinct
subproblem -- is the opening idea of the next chapter, where it is
applied systematically across dozens of problems, progressing from this
kind of memoized top-down recursion to fully iterative tabulation and,
where possible, space-optimized tabulation.

\subsection*{Algorithm}

\begin{enumerate}
  \item $\textsc{Solve}(\textit{start})$:
  \begin{enumerate}
    \item If $\textit{start}=|s|$: return \texttt{true}.
    \item For $\textit{end}$ from $\textit{start}$ to $|s|-1$: if
          $s[\textit{start}..\textit{end}] \in \textit{wordDict}$ and
          $\textsc{Solve}(\textit{end}+1)$: return \texttt{true}.
    \item Return \texttt{false}.
  \end{enumerate}
  \item Call $\textsc{Solve}(0)$.
\end{enumerate}

\subsection*{Dry run}

\dryrun{Take $s = \texttt{"leetcode"}$, $\textit{wordDict} =
\{\texttt{"leet"}, \texttt{"code"}\}$.}

\begin{itemize}
  \item $\textsc{Solve}(0)$: try $\textit{end}=0..3$:
        $s[0..3]=\texttt{"leet"}$ is in the dictionary. Check
        $\textsc{Solve}(4)$.
  \item $\textsc{Solve}(4)$: try $\textit{end}=4..7$:
        $s[4..7]=\texttt{"code"}$ is in the dictionary. Check
        $\textsc{Solve}(8)$.
  \item $\textsc{Solve}(8)$: $\textit{start}=8=|s|$. Return
        \texttt{true}.
  \item This propagates straight back: $\textsc{Solve}(4)$ returns
        \texttt{true}; $\textsc{Solve}(0)$ returns \texttt{true}.
\end{itemize}

The string segments as \texttt{"leet"} + \texttt{"code"}.

\subsection*{C++ implementation}

\begin{lstlisting}
#include <bits/stdc++.h>
using namespace std;

bool solve(int start, string& s, unordered_set<string>& wordSet) {
    if (start == (int)s.size()) return true;

    for (int end = start; end < (int)s.size(); end++) {
        string piece = s.substr(start, end - start + 1);
        if (wordSet.count(piece) && solve(end + 1, s, wordSet)) {
            return true;
        }
    }
    return false;
}

bool wordBreak(string s, vector<string>& wordDict) {
    unordered_set<string> wordSet(wordDict.begin(), wordDict.end());
    return solve(0, s, wordSet);
}

int main() {
    vector<string> wordDict = {"leet", "code"};
    cout << (wordBreak("leetcode", wordDict) ? "true" : "false") << endl;
    return 0;
}
\end{lstlisting}

\begin{complexitybox}
\textbf{Time Complexity.} $O(2^n)$ in the worst case for this naive
(unmemoized) version, due to the overlapping-subproblem blowup described
above. With memoization keyed by $\textit{start}$ (an array or hash map
of size $O(n)$, each computed once in $O(n)$ work to try all
$\textit{end}$ positions), this drops to $O(n^2)$.
\textbf{Space Complexity.} $O(n)$ for the recursion depth, plus
$O(n)$ for the memoization cache once added.
\end{complexitybox}

\begin{takeawaybox}
This closing problem of the Recursion chapter is a deliberate hinge:
everything about its recursive structure is already correct and familiar
from this chapter's earlier sections, but its overlapping subproblems
expose recursion's biggest weakness -- redundant recomputation -- setting
up exactly the problem that memoization and tabulation, covered next, are
designed to solve. Carry this recursion forward unchanged; only the
bookkeeping around it is about to be upgraded.
\end{takeawaybox}
